\section{Adaptive Gating Policy}
\label{sec:method}
% =============================================================================

The gating policy is a lightweight network trained on top of a frozen
AlphaZero planner that selects, at each meta-step, a planning budget
$K \in \Kset$.
We described the variable-delay setting it operates in formally in
Section~\ref{sec:envs}; here we describe the policy itself---what it
observes, how it is structured, and how it is trained.

Figure~\ref{fig:scaling} shows the tradeoff the gating policy must
manage: additional MCTS simulations usually improve action quality, but
they also take longer to run.
The planner already handles action selection; the gate only decides,
state by state, how long to run it.
We therefore train the gate on top of a frozen planner rather than
jointly optimizing both components.

\subsection{Observations and architecture}

At each meta-step the gating policy takes three inputs: (1) the raw
game grid, (2) the planner's intermediate spatial features, extracted
from its frozen trunk and used as a compact summary of what the planner
currently infers about the position, and (3) the planner's scalar value
estimate $V(s_t)$.
In clock environments the observation additionally includes the
remaining time budget for both players.
Together these inputs expose the raw state, the planner's internal
representation, and, where applicable, the current time pressure.
We do not feed the planner's policy logits to the gate, so the budget
decision depends on state and planner features rather than directly on
the planner's preferred action.

A lightweight network processes these inputs and produces a categorical
distribution over $K$ together with a scalar baseline value used during
training.
Architecture details, including the recurrent variant used for Speed
Hex, are in Appendix~\ref{app:implementation}.

\subsection{Training}

Training proceeds in two phases.
We first train AlphaZero-style base planners for each environment using
standard expert iteration (Appendix~\ref{app:az_primer}) and select
frozen base checkpoints by cross-evaluation across delay budgets
$K \in \Kset$ (Appendix~\ref{app:cross_eval}).
We then freeze the selected planner and train the gating policy with
PPO~\citep{schulman2017proximal} on top of it.

Rollouts are collected over full episodes. At each meta-step the gate
samples $K_t \sim \pigate(\cdot \mid s_t)$,
the agent executes the resulting option---$K_t{-}1$ filler actions
followed by the planner's chosen action---and the meta-reward $R_t$
defined in Section~\ref{sec:envs} is accumulated.
Variable-duration advantages with per-meta-step discount $\gamma^{K_t}$
are computed before each round of PPO updates, adapting Generalized
Advantage Estimation~\citep{schulman2015high} to the SMDP introduced in
equation~(\ref{eq:smdp}).
Only the gating network is updated; no gradient flows to the planner.

To keep training efficient, we evaluate all $|\Kset|$ option
transitions in parallel at each meta-step.
This preserves static shapes and keeps the training loop easy to
compile.
The full JAX implementation---including the use of \texttt{lax.scan}
for both MCTS inner loops and PPO rollouts, \texttt{vmap} for
environment parallelism, and \texttt{pmap} for device parallelism---is
described in Appendix~\ref{app:implementation}.
